\documentclass[11pt]{article}

\usepackage[margin=1in]{geometry}
\usepackage{array}
\usepackage{amsmath,amssymb,amsfonts}
\usepackage{booktabs}
\usepackage{algorithm}
\usepackage{algorithmic}
\usepackage{enumitem}
\usepackage{hyperref}
\usepackage{xcolor}

\hypersetup{colorlinks=true, linkcolor=blue!45!black, citecolor=blue!45!black}

\newcommand{\Dold}{\mathcal{D}_{\mathrm{old}}}
\newcommand{\Dnew}{\mathcal{D}_{\mathrm{new}}}
\newcommand{\X}{\mathbf{X}}
\newcommand{\z}{\mathbf{z}}
\newcommand{\VIMP}{\mathrm{VIMP}}
\newcommand{\MMD}{\mathrm{MMD}}
\newcommand{\CS}{\textsc{cs}}
\newcommand{\CD}{\textsc{cd}}

\title{\bf FSDS-SoT: A Complete Write-Up\\[0.4em]
\large Actionable Topology, Critical-Path Budgeting, and Incremental Value}
\author{Prototype documentation}
\date{}

\begin{document}
\maketitle

\begin{abstract}
We document a production-oriented landing of Feature Selection for Distribution
Shift (FSDS) on Skeleton-of-Thought (SoT). Two objects carry essentially all
actionable value: \textbf{(1)~grouping granularity}---a suggested parallel width
$K$, merge rules, and ranked execution order within each cluster from branch
embeddings; and \textbf{(2)~a critical-path budget executor}---per-branch
expansion tokens $L_b$, model tier, and check budget from shift and quality.
Batch covariate monitoring, merge at assembly time, and overlap-based fallback
close the loop but do not replace these two. We provide summary tables, dashboard
semantics, incremental-value categories beyond cost reduction, and pointers to
the reference implementation (\texttt{Python/fsds\_sot/}) and agentic demo.
Meta-learners are used for distance estimation only, not causal identification.
\end{abstract}

\section{Scope and two core objects}
\label{sec:scope}

SoT generates a skeleton $\{s_1,\ldots,s_B\}$ and expands branches, often in
parallel. FSDS does not replace the LLM; it schedules \emph{when} to parallelise,
\emph{how many} independent groups to run, and \emph{how much} compute each
branch receives. Business rules (minimum/maximum width, latency SLA, compliance)
clip data-driven suggestions:
\begin{equation}
  K_{\mathrm{final}} = \mathrm{clip}\big(K_{\mathrm{FSDS}},\, K_{\min}^{\mathrm{biz}},\, K_{\max}^{\mathrm{biz}}\big).
\end{equation}

\paragraph{Object~1 (topology).} Embedding coupling $\Rightarrow$ suggested $K$,
cluster membership, merge pairs, sequential order inside clusters, parallel gate,
straggler flags.

\paragraph{Object~2 (resources).} Shift $s_b$ and quality $q_b$ $\Rightarrow$
$L_b$, tier, checks, optional zero-sum reclaim under a total token budget.

\section{Procedure stack (five steps)}
\label{sec:procedures}

\begin{enumerate}[leftmargin=*]
  \item[\textbf{(0)}] \textbf{Batch covariate monitoring} --- $\Dold$ vs.\
  $\Dnew$ on trace features $\phi(\mathcal{S})$; RF-domain $\VIMP$, $\MMD$-LOCO
  (no outcome $Y$ required).
  \item[\textbf{(1)}] \textbf{Grouping granularity} --- Object~1 on
  $\{\z_b\}_{b=1}^B$.
  \item[\textbf{(2)}] \textbf{Critical-path budget executor} --- Object~2 after
  topology is fixed.
  \item[\textbf{(3)}] \textbf{Merge / aggregation} --- deduplicate and reconcile
  parallel expansions (GoT-style at merge time).
  \item[\textbf{(4)}] \textbf{Confidence and fallback} --- overlap, bootstrap
  CI, stability; abstain or covariate-only when unreliable.
\end{enumerate}

\section{Actionable outputs at a glance (tables)}
\label{sec:action-tables}

Table~\ref{tab:obj1-topology} summarises the primary topology object: suggested
parallel width $K$, merge rules, and execution order inside each cluster.
Table~\ref{tab:obj2-budget} gives the budget executor outputs.
Table~\ref{tab:incremental-value} lists incremental value beyond cost reduction.
Table~\ref{tab:dashboard-panels} maps the four-panel dashboard to decisions.

\begin{table}[t]
\centering
\small
\caption{Object~1 --- grouping granularity (embedding topology). $K$ is a
\emph{suggestion}; $K_{\mathrm{final}}=\mathrm{clip}(K_{\mathrm{FSDS}},K_{\min}^{\mathrm{biz}},K_{\max}^{\mathrm{biz}})$.}
\label{tab:obj1-topology}
\begin{tabular}{@{}p{0.22\linewidth}p{0.30\linewidth}p{0.38\linewidth}@{}}
\toprule
Output & How it is computed & Actionable insight \\
\midrule
Suggested size $K$ &
Coupling graph on $\{\z_b\}$; connected components / spectral clusters &
How many \emph{parallel groups} to run; clip by business min/max width \\
Cluster membership &
Cluster labels per skeleton point $b$ &
Which branches share a parallel slot vs.\ must not race \\
Merge pairs &
High coupling + concept-redundant ($C_{ij}\le\varepsilon_c$) &
Collapse duplicate skeleton points before expand (no double-counting) \\
Sequential order &
Topological / dependency order \emph{within} a cluster &
Ranked expand sequence when intra-cluster coupling is high \\
Parallel vs.\ sequential &
$P=1-\overline{H}_{ij}$, redundancy $R$, balance $\mathrm{Bal}$ &
Allow SoT parallel only if gate passes; else CoT or cluster-internal chain \\
Straggler flag &
Imbalanced predicted shift or length across $b$ &
Split or cap a branch before it dominates critical path \\
\bottomrule
\end{tabular}
\end{table}

\begin{table}[t]
\centering
\small
\caption{Object~2 --- critical-path budget executor (after topology is fixed).}
\label{tab:obj2-budget}
\begin{tabular}{@{}p{0.22\linewidth}p{0.30\linewidth}p{0.38\linewidth}@{}}
\toprule
Output & Driver & Actionable insight \\
\midrule
Expansion tokens $L_b$ &
Normalized shift $s_b$, quality $q_b$; clip to $C_{\mathrm{lat}}$ &
Set \texttt{max\_new\_tokens} per branch; longest $L_b$ sets span \\
Model tier &
High $s_b$ and risk $s_b(1-q_b)$ &
Route branch to small / medium / large model \\
Check budget &
Risk $\propto s_b(1-q_b)$; minimum $1$ if low shift but low $q_b$ &
Targeted verify / tool-validity calls (not uniform on every branch) \\
Zero-sum reclaim &
Low-shift, high-$q_b$ branches &
Reallocate saved tokens to drifting branches under fixed total budget \\
\bottomrule
\end{tabular}
\end{table}

\begin{table}[t]
\centering
\small
\caption{Incremental value of FSDS-SoT vs.\ uniform SoT (justify with A/B or
frontier: success $\times$ cost $\times$ span).}
\label{tab:incremental-value}
\begin{tabular}{@{}llp{0.42\linewidth}@{}}
\toprule
Category & Primary object & What improves (beyond lower token bill) \\
\midrule
Latency / SLA & Object~2 &
Shorter critical path (span); fewer straggler timeouts under parallel decode \\
Reliability under drift & Object~2 (+ $Y$) &
Higher success / tool-validity when retrieve--act branches shift \\
Observability & Batch $\VIMP$ + Object~1 &
Localised drift tickets; faster MTTR vs.\ aggregate accuracy alone \\
Safety / compliance & Object~2 &
Directed checks on high-risk cells; minimum guardrail on LoS$\cdot$LoQ \\
Capacity & Object~1 + Object~2 &
Fewer bad parallel merges; more effective QPS per GPU slot \\
Strategy safety & Object~1 &
Fallback to sequential when decomposability fails (avoid SoT tail risk) \\
Learning loop (slow) & Batch monitoring &
Targeted retrain / AGOD on persistent drift slices \\
\bottomrule
\end{tabular}
\end{table}

\begin{table}[t]
\centering
\small
\caption{Dashboard panels: embedding-only vs.\ mixed signals.}
\label{tab:dashboard-panels}
\begin{tabular}{@{}llll@{}}
\toprule
Panel & Data source & Suggestion type & Typical action \\
\midrule
1 Coupling heatmap & Branch embedding & $K$, merge, sequentialise &
Re-plan groups before expand \\
2 PCA clusters & Branch embedding & Same partition as Panel~1 &
Visual audit of cluster assignment \\
3 Budget bars & Embedding shift + quality & $L_b$, tier, checks &
Deploy per-branch scheduler limits \\
4 Batch VIMP & Trace features $\Dold$ vs.\ $\Dnew$ & Default policy drift &
Adjust templates, caps, monitoring cadence \\
\bottomrule
\end{tabular}
\end{table}

\section{Decomposability and suggested $K$}
\label{sec:decomposability}

On branch embeddings, form $H_{ij}=\widehat{\mathrm{HSIC}}_{\mathrm{norm}}(\z_i,\z_j)$,
threshold at $\delta$, and count connected components $K$. Decomposability degree
$\mathcal{D}=K/B$. Gate on redundancy $R$, parallelizability $P=1-\overline{H}_{ij}$,
and balance $\mathrm{Bal}$ before allowing full parallel SoT.

\section{Critical-path budget (formulas)}
\label{sec:budget-formulas}

Parallel latency is governed by span, not total tokens:
\begin{equation}
  T_{\infty} = t_{\mathrm{skel}} + \max_{c\in\Pi}\sum_{b\in\mathrm{chain}(c)} L_b\, t_{\mathrm{tok}} + t_{\mathrm{merge}}.
\end{equation}
Allocate
\begin{equation}
  L_b = \mathrm{clip}(\kappa\, s_b\, q_b,\, L_{\min},\, C_{\mathrm{lat}}),
  \qquad
  \mathrm{risk}_b \propto s_b(1-q_b),
\end{equation}
with check budget clipped to $[1,5]$ and tier from joint high shift and risk.
LoS$\cdot$LoQ (low shift, low quality): retain at least one check; do not upscale
the model without evidence of new content.

\section{Incremental value: how to justify}
\label{sec:justify}

\paragraph{Counterfactual.} Compare FSDS-SoT to uniform SoT at fixed model and
skeleton template; attribute lift to Objects~1--2 via ablation (disable gate or
disable budget executor).

\paragraph{Evidence ladder.} (i)~Mechanism demo on agentic traces
(\texttt{demo\_fsds\_sot\_agentic\_incremental.py}); (ii)~batch $\VIMP$ aligned
with retrieve/act drift; (iii)~production A/B on success, span, cost.

\paragraph{Agentic demo (sketch).} Skeleton branches
\texttt{plan, retrieve, reason, act, verify}; live batch injects embedding drift
on tool branches. Typical demo direction: large cost/check savings with span and
success validated on real $Y$ in deployment.

\section{SoT usage: efficiency beyond compute cost and utilization}
\label{sec:sot-efficiency}

Cost and GPU utilisation answer ``how much we pay'' and ``how full the cluster
is.'' \emph{Efficiency} answers ``how much useful reasoning we get per unit of
wall-clock and per unit of generation''---the operational reason to adopt SoT in
the first place, amplified when FSDS schedules topology and budgets. We separate
four efficiency notions that SoT practitioners already care about; FSDS makes
each one measurable and steerable.

\paragraph{Latency efficiency (span).} Under parallel decode, user-visible latency
is dominated by the critical path $T_\infty$ (Section~\ref{sec:budget-formulas}),
not by total tokens $T_1$. Uniform SoT often assigns every branch the same
expansion cap, so a single straggler branch fixes $T_\infty$ even when other
branches finish early---parallelism does not help the SLA. FSDS shortens or caps
branches with low shift, splits or reorders high-coupling pairs (Object~1), and
concentrates length on branches that actually carry new information (Object~2).
With $P$ parallel slots, wall-clock satisfies Brent's bound
$T_P \le T_\infty + (T_1-T_\infty)/P$; reducing $T_\infty$ is the direct lever on
p99 and timeout rate \emph{without} requiring higher utilisation.

\paragraph{Parallel efficiency ($K$ vs.\ $P$).} SoT's benefit assumes the skeleton
decomposes into $K$ weakly coupled clusters that can run concurrently. When
$K_{\mathrm{final}}>P$, execution proceeds in waves; idle slots and merge
contention are pure efficiency losses. Object~1 estimates $K$, flags merge pairs,
and enforces cluster-internal order so that each wave runs \emph{meaningful}
parallel work rather than racing branches that must later be reconciled. This is
efficiency of \emph{parallel form}: fewer false starts, fewer merge retries, better
slot occupancy in each wave.

\paragraph{Generation efficiency (quality per token).} Efficiency is not only
``fewer tokens'' but ``more outcome per token.'' Aligning $L_b$ and tier with
shift~$\times$~quality puts capacity on branches whose embeddings indicate new
content or low verifier confidence, and makes terse expansions on stable,
high-quality branches---raising success or tool-validity per generated token
rather than uniformly inflating or deflating all branches. Targeted checks
(Object~2) add verification only on high-risk cells, improving \emph{verification
efficiency} (fewer verify calls per successful episode).

\paragraph{Workflow efficiency (time-to-action).} Before expansion, the dashboard
turns embedding topology into explicit decisions: suggested $K$, merge list,
ranked order within clusters, and per-branch $(L_b,\mathrm{tier},\mathrm{checks})$.
Operators and schedulers skip trial-and-error uniform caps and post-hoc firefighting;
batch $\VIMP$ (Panel~4) separates ``change default skeleton policy'' from
``change this episode's branch budget,'' reducing mean time to diagnose drift.
That is efficiency of the \emph{control loop}, not of the GPU kernel.

\begin{table}[t]
\centering
\small
\caption{Efficiency dimensions for FSDS-scheduled SoT (orthogonal to raw cost and
utilisation).}
\label{tab:sot-efficiency}
\begin{tabular}{@{}p{0.24\linewidth}p{0.34\linewidth}p{0.34\linewidth}@{}}
\toprule
Efficiency notion & What SoT alone assumes & What FSDS adds \\
\midrule
Latency (span) &
Equal branch lengths; critical path = longest branch &
Cap/straggler control; shift-weighted $L_b$; cluster ordering \\
Parallel form &
All skeleton points expand in parallel &
Suggested $K$, merge, sequentialise; gate when $\mathcal{D}$ is low \\
Quality / token &
Same model and length everywhere &
Tier routing; risk-based checks; terse stable branches \\
Control loop &
Manual tuning of width and caps &
Tables~\ref{tab:obj1-topology}--\ref{tab:dashboard-panels}; drift-local actions \\
\bottomrule
\end{tabular}
\end{table}

\paragraph{Usage takeaway.} Use FSDS-SoT when SoT is already in the stack but
episodes show uneven branch shift, tool-branch drift, straggler timeouts, or
unexplained merge/rework: Object~1 fixes \emph{how many groups and in what order};
Object~2 fixes \emph{how much generation and verification each group receives}.
Efficiency gains appear as shorter critical paths, higher valid parallel width,
better outcome per token, and faster operational response---complementary to,
not redundant with, compute-cost savings already documented elsewhere.

\paragraph{In one paragraph (SoT usage).}
Skeleton-of-Thought wins when branches are weakly coupled and expansion budgets
match where the answer actually lives; it loses when every branch gets the same
length, the longest branch sets the SLA, or parallel runs race and must be merged
again. FSDS-SoT reads branch embeddings once, suggests how many parallel groups
$K$ to run (merge and order within each group), then assigns each group its own
expansion cap, model tier, and check count from shift and quality---so wall-clock
tracks the critical path, parallel slots run meaningful work, and each generated
token is more likely to move the task forward. That is the everyday usage case:
SoT for speed, FSDS for \emph{where} to parallelise and \emph{how much} to spend
on each skeleton point.

\section{ToT vs.\ SoT (one row each)}
\label{sec:tot-sot}

\begin{table}[h]
\centering
\small
\begin{tabular}{lll}
\toprule
 & SoT & ToT \\
\midrule
Width & Suggested $K$ / branches & Branching factor per node \\
Depth & Per-branch $L_b$ (tokens) & Search depth / prune policy \\
Extra & --- & Re-calibrate value head if $\CD$ on value features \\
\bottomrule
\end{tabular}
\end{table}

\section{Implementation and artifacts}
\label{sec:impl}

\begin{itemize}[leftmargin=*]
  \item Code: \texttt{Python/fsds\_sot/} (pipeline, budget, viz, agentic DGP).
  \item Demos: \texttt{demo\_fsds\_sot.py}, \texttt{demo\_fsds\_sot\_viz.py},
  \texttt{demo\_fsds\_sot\_agentic\_incremental.py}.
  \item Docs: \texttt{docs/FSDS\_SoT\_LANDING.md},
  \texttt{docs/FSDS\_SoT\_AGENTIC\_INCREMENTAL\_VALUE.md}.
  \item Extended manuscript: \texttt{FSDS\_SoT\_ToT\_manuscript.tex} with appendix
  \texttt{FSDS\_SoT\_ToT\_supplement.tex}.
  \item Agent applications \& closed loop: \texttt{FSDS\_agent\_reasoning\_applications.tex};
  demo \texttt{demo\_fsds\_closed\_loop.py}.
\end{itemize}

% FSDS on agent reasoning: extended applications and closed-loop self-iteration
% Compile standalone or \input into FSDS_SoT_complete_writeup.tex after macros.

\section{Agent-reasoning applications beyond SoT scheduling}
\label{sec:agent-apps}

FSDS treats each agent episode as a structured feature object and compares a
reference batch to a live batch. Skeleton-of-Thought (SoT) is one instance:
branch embeddings supply Object~1 (topology) and Object~2 (budget). The same
pattern applies wherever an agent exposes \emph{segmentable} states---steps,
tool observations, sub-agents, or memory windows---and a batch label
(old vs.\ new traffic, model version, or policy regime).

\begin{table}[t]
\centering
\small
\caption{Agent-reasoning application map. Each row uses FSDS for monitoring
(covariate $\VIMP$) and, when $Y$ is defined, concept attribution; scheduling
actions follow the separation principle ($\CS$ $\to$ structure/route;
$\CD$ $\to$ content/scoring/verify).}
\label{tab:agent-apps}
\begin{tabular}{@{}p{0.18\linewidth}p{0.26\linewidth}p{0.22\linewidth}p{0.26\linewidth}@{}}
\toprule
Pattern & Natural batch / segments & Proxy $Y$ & FSDS-driven action \\
\midrule
SoT / plan--execute &
Skeleton branches $\{\z_b\}$ &
Verifier, task success &
Suggested $K$, merge/order; $L_b$, tier, checks \\
ReAct / tool loop &
Each \texttt{Observation} span &
Tool validity, grounding &
Continue vs.\ re-call tool; block noisy obs \\
Multi-agent &
Per-agent message embeddings &
Handoff success, subtask $Y$ &
Which agent to activate; merge conflicts \\
Self-consistency &
$N$ sample chains &
Vote margin, dispersion &
Adaptive $N$; targeted resampling \\
Reflexion &
Draft vs.\ critique embeddings &
Self-critic score &
Iterate again? rewrite which span \\
RAG agent &
Per retrieved chunk / sub-query &
Citation support &
Re-retrieve depth; drop bad chunks \\
Supervisor--worker &
Worker outputs at step $t$ &
Supervisor approval &
Re-route worker; adjust parallel width \\
Streaming memory &
Time windows of embeddings &
PRM step score &
KV invalidate; re-encode modality \\
Human escalation &
High-risk cells (Panel~3) &
Escalation outcome &
Abstain / human-in-the-loop trigger \\
Speculative drafts &
Draft vs.\ verify embeddings &
Accept/reject draft &
Which draft to expand; early reject \\
Code agent &
Edit spans / test logs &
Tests pass, lint &
Re-run tests; patch which file span \\
Tool registry &
Call features vs.\ schema version &
Parse/422 errors &
Route to fallback API; freeze tool \\
Multi-hop RAG &
Hop-$k$ chunk embeddings &
Answer groundedness &
Stop hops; widen query at hop $k^\star$ \\
\bottomrule
\end{tabular}
\end{table}

\paragraph{Opportunity criteria.} A new application is a good fit when
(i)~segments are cheap to embed or featurise online,
(ii)~a reference batch exists (rolling window, canary, or A/B control),
(iii)~an action lever exists (budget, route, verify, retrieve, freeze), and
(iv)~overlap between batches is sufficient or covariate-only mode is acceptable.

\section{Closed-loop self-iteration}
\label{sec:closed-loop}

Production use is not a one-shot attribution. We close the loop: \emph{measure
$\to$ attribute $\to$ act $\to$ re-measure $\to$ escalate or accept}. This is the
self-iteration protocol that ties batch monitoring (procedure~0), SoT topology
(procedure~1), budget execution (procedure~2), merge (procedure~3), and fallback
(procedure~4) into a single control cycle.

\begin{algorithm}[t]
\caption{FSDS closed-loop iteration (per traffic window)}
\label{alg:closed-loop}
\begin{algorithmic}[1]
\STATE \textbf{Input:} reference window $\Dold$, live window $\Dnew$, policy $\pi_0$
\STATE Extract $\X$, optional $Y$; fit covariate attribution; test overlap
\IF{overlap poor}
  \STATE Report covariate-only; $\pi \leftarrow \pi_{\mathrm{safe}}$; \textbf{return}
\ENDIF
\STATE Localise drift shares on segments (nodes, branches, agents, windows)
\STATE Select intervention $a \sim \pi(\cdot \mid g)$ from cost ladder
\STATE Apply $a$ (re-plan, budget, route, retrieve, verify, rollback)
\STATE Collect $\Dnew'$ on the next window; re-run FSDS($\Dold$, $\Dnew'$)
\IF{target drift share $\downarrow$ and KPI non-inferior}
  \STATE \textbf{Accept} $a$; set $\Dold \leftarrow$ decay($\Dold$, $\Dnew'$)
\ELSE
  \STATE \textbf{Escalate} one rung on ladder or invoke fallback
\ENDIF
\STATE Persist $(g, a, \Delta\mathrm{KPI})$ for AGOD / offline retrain (slow loop)
\end{algorithmic}
\end{algorithm}

\paragraph{Intervention ladder (self-iteration order).}
\begin{enumerate}[leftmargin=*]
  \item \textbf{Reallocate} --- $L_b$, tier, checks, retrieval depth (cheapest).
  \item \textbf{Re-topology} --- merge, reorder, change $K_{\mathrm{final}}$, disable parallel.
  \item \textbf{Re-ground} --- re-retrieve, tool re-call, verify/backtrack span.
  \item \textbf{Re-configure} --- prompt/template, rollback model version.
  \item \textbf{Re-learn} --- AGOD adapter update, targeted data collection.
\end{enumerate}
Each iteration consumes one window; escalation stops when overlap fails or the
ladder is exhausted (abstain / human).

\paragraph{Self-iteration metrics.} Track (i)~drift share on the targeted feature
group, (ii)~KPI vs.\ pre-intervention baseline (success, span, cost),
(i)~intervention acceptance rate, (iv)~mean rungs to stabilise. A stable system
shows decaying drift share without KPI regression---evidence that the loop is
converging, not chasing noise.

\section{Implemented self-iteration stack (code)}
\label{sec:impl-loop}

\begin{table}[t]
\centering
\small
\caption{Repository modules for closed-loop agent reasoning (branch
\texttt{cursor/fsds-sot-tot-manuscript-7451}).}
\label{tab:impl-loop}
\begin{tabular}{@{}ll@{}}
\toprule
Module / demo & Role in the loop \\
\midrule
\texttt{closed\_loop.py} &
\texttt{iterate\_once}, \texttt{run\_self\_iteration}, ladder, reference decay \\
\texttt{applications.py} &
ReAct, plan--execute, multi-agent, RAG hops, self-consistency segments \\
\texttt{demo\_fsds\_closed\_loop.py} & SoT batch: two-window drift decay \\
\texttt{demo\_fsds\_react\_agent.py} & ReAct segment budgets + economics \\
\texttt{demo\_fsds\_rag\_multihop\_loop.py} &
Multi-hop reground until \texttt{stabilize\_threshold} \\
\texttt{demo\_fsds\_self\_consistency\_loop.py} &
Dispersion $\downarrow$ $\Rightarrow$ adaptive $N^\star$ + loop accept \\
\bottomrule
\end{tabular}
\end{table}

\paragraph{Reference update.} After an accepted window, $\Dold \leftarrow
\alpha \Dold + (1-\alpha)\Dnew'$ with $\alpha\in[0.8,0.95]$ so the control
system tracks a moving baseline without forgetting recent fixes.

\section{Roadmap: remaining hooks}
\label{sec:roadmap}

Production hooks still to wire (same estimators, new $\phi$ only):
\begin{itemize}[leftmargin=*]
  \item \textbf{Live KPI guard} --- non-inferiority test on conversion / resolution
  before accepting an intervention.
  \item \textbf{Code-agent span PO-risk} --- edit-hunk $Y$ from CI logs.
  \item \textbf{Tool-registry batch} --- schema-version covariate block in Panel~4.
\end{itemize}
SoT remains the reference implementation; other rows in
Table~\ref{tab:agent-apps} reuse the same estimators and differ only in
$\phi(\cdot)$ and $\pi$.

% One-page boss brief — compile standalone or \input from complete writeup
\section*{Stakeholder brief: six points (FSDS agent SoT POC)}
\small
\textbf{Run numbers:} \texttt{python3 demo\_fsds\_boss\_six\_points.py} $\rightarrow$
\texttt{artifacts/boss\_delivery\_six\_points.json}.

\paragraph{1. Problem \& two objects.}
Live agent traces drift vs reference; uniform SoT wastes tokens/checks.
\textbf{Object~1:} topology gate ($K$, merge/order).
\textbf{Object~2:} critical-path budget ($L_b$, tier, checks) from shift$\times$quality$\times$need.

\paragraph{2. Critical-path budget.}
$L_b=\mathrm{clip}(\kappa s_b q_b \,\mathrm{need}_b, L_{\min}, C_{\mathrm{lat}})$;
parallel latency $T_\infty=\max_b L_b+\cdots$.
Savings often appear in \emph{total work} $\sum_b L_b$, not span, when one branch saturates $C_{\mathrm{lat}}$.

\paragraph{3. $\sim$28\% total tokens.}
Water-filling + global budget cap shrinks low-risk branches; eval mean
$2600\rightarrow 1870$ tokens/ep ($\approx 28\%$ on synthetic agentic batch).

\paragraph{4. Checks halved.}
Uniform: $2$ per branch ($10$/ep). FSDS: $\mathrm{round}(1+2s_b(1-q_b))$;
eval mean $\approx 5$ checks/ep.

\paragraph{5. Pareto + MCTS + entropy.}
Entropy-regularized weights
$w=(1-\lambda)\mathrm{softmax}(\log raw/\tau)+\lambda/B$;
MCTS/local search on $(\kappa,\lambda,\mathrm{floor})$ for success-rate
(\texttt{mcts\_search.py}).
\paragraph{5b. ROI.}
\texttt{net\_economic\_gain} = $\Delta(\mathrm{success}\cdot v_{\mathrm{succ}})$ +
variable cost savings $-$ overhead; unified tier/check pricing (\texttt{pricing.py}).

\paragraph{6. Delivery.}
Code: \texttt{fsds\_sot/}; PR \#12; next: wire real resolution/tool\_ok into guard, canary adopt.

% Online PFI delivery section — \input from FSDS_SoT_complete_writeup.tex
\section{Online PFI: streaming covariate attribution}
\label{sec:online-pfi}

\paragraph{Logic stack.}
Each agent episode yields trace features $\X=\phi(\mathcal{S})$.
On window $t=1,2,\ldots$, compare live batch $\Dnew$ to reference $\Dold$
(optional exponential decay after each step).
The covariate plane fits a domain RF $W\sim \X$ and computes
\textbf{online permutation feature importance} $\VIMP_t$ (same estimator as
batch Panel~4, repeated over time).
Parallel signals: $\MMD^2$, domain AUC, overlap ESS, and an
\textbf{onlineRFPerm} classifier two-sample $p$-value from label permutations.

\begin{algorithm}[t]
\caption{Online PFI stream (implementation: \texttt{online\_pfi.py})}
\label{alg:online-pfi}
\begin{algorithmic}[1]
\STATE Input: $\X_{\mathrm{ref}}$, live windows $\{\X^{(t)}\}$, decay $\alpha$
\FOR{$t=0,1,\ldots$}
  \STATE $\VIMP_t, \mathrm{AUC}_t, \MMD^2_t \leftarrow \textsc{CovariateAttribution}(\X_{\mathrm{ref}}, \X^{(t)})$
  \STATE $p_t \leftarrow \textsc{RFPermPvalue}(\X_{\mathrm{ref}} \cup \X^{(t)})$
  \IF{$p_t < \alpha_{\mathrm{trig}}$}
    \STATE emit drift signal; route to closed-loop ladder \S\ref{sec:closed-loop}
  \ENDIF
  \STATE $\X_{\mathrm{ref}} \leftarrow \alpha \X_{\mathrm{ref}} + (1-\alpha)\X^{(t)}$
\ENDFOR
\end{algorithmic}
\end{algorithm}

\paragraph{Visualization.}
Four-panel dashboard: AUC and $\MMD^2$ vs.\ window; $p$-value stream;
heatmap of top-$k$ $\VIMP$ over time
(\texttt{plot\_online\_pfi.py}, artifact \texttt{07\_online\_pfi\_dashboard.png}).

\paragraph{Delivery package.}
Run \texttt{package\_online\_pfi\_delivery.py}; outputs under
\texttt{delivery/online\_pfi/}. See \texttt{docs/FSDS\_ONLINE\_PFI\_DELIVERY\_PACKAGE.md}.


\section*{Acknowledgement of scope}
FSDS attributes nonstationarity in features and embeddings used for scheduling;
it is not next-token gradient interpretability. Tables~\ref{tab:obj1-topology}--\ref{tab:dashboard-panels}
are normative for operations; empirical numbers in demos are illustrative until
replaced by production A/B.

\end{document}
